\chapter{Master Integrals}\label{ap:masterIntegrals}

This appendix collects the definitions and the analytic evaluation of the master integrals that appear
in the integrated antenna functions of Chapter~\ref{ch:workedexample}, organised by antenna family:
the two-loop $A$-type integrals of the $A_2^2$ set, the one-loop phase-space $V$-type integrals of the
$A_3^1$ set, and the four-particle phase-space $R$-type integrals of the four-parton antennae. The
only massless master integral treated in the main text is $R_3$, which is needed for the worked
example of Subsection~\ref{sec:matDerivation}, and the massive master integrals are presented together
with the massive $A_3^0$ antenna in Section~\ref{sec:massiveA30}. The notation for the phase-space
measures and for the normalisation factors $S_\Gamma$ and $S_{\Gamma,2}$ is collected in
Section~\ref{sec:MInotation}, and the scalar one-loop functions $B_0$, $C_0$ and $D_0$
are discussed in Appendix~\ref{subsec:PaVe}. The derivation of $R_{8,a}$, the most involved of the
massless master integrals, is given in full detail in Section~\ref{sec:MIfourParton} as an example of
the method used for the others.

\begin{table}[h!]
\footnotesize
\centering
\caption{Overview of the master integrals used in the build- and integration-stage by each antenna type \cite{Gehrmann-DeRidder:2005btv,GehrmannDeRidder0311276}. The entries marked $^\dagger$ arise only from the $A_2^1$ ultraviolet counterterm attached at the public build boundary; the two-loop source itself is reduced by IBP to the $A$-type masters.}
\renewcommand{\arraystretch}{1.3}
  \begin{tabular}{|c|c|c|} 
   \hline
   \textbf{Antenna type(s)} & \textbf{Build-Side MI} & \textbf{Integration-Side MI} \\ 
   \hline
   $A_2^0$ & -- & -- \\ \hline
   $A_3^0$ & -- & $R_3$ \\ \hline
   $A_4^0$, $\tilde A_4^0$, $B_4^0$, $C_4^0$ & -- & $R_4$, $R_6$, $R_{8,a}$, $R_{8,b}$ \\ \hline
   $A_2^1$ & -- & $B_0$, $C_0$ \\ \hline
   $A_3^1$, $\tilde A_3^1$, $\hat A_3^1$ & $B_0$, $C_0$, $D_0$ & $V_{5,a}$, $V_{5,b}$, $V_8$ \\ \hline
   $A_2^2$ & $A_{22,\text{LO}}$, $A_3$, $A_4$, $B_0^\dagger$, $C_0^\dagger$ & $A_{22,\text{LO}}$, $A_3$, $A_4$ \\ \hline
   $\tilde A_2^2$ & $A_{22,\text{LO}}$, $A_3$, $A_4$, $A_6$ & $A_{22,\text{LO}}$, $A_3$, $A_4, A_6$ \\ \hline
   $\hat A_2^2$ & $A_4$, $B_0^\dagger$, $C_0^\dagger$ & $A_4$ \\ \hline
   $\breve A_2^2$ & $A_{22,\text{LO}}$ & $A_{22,\text{LO}}$ \\ \hline
  \end{tabular}
   
\label{tab:MI}
\end{table}

\section{Phase-Space Notation and Normalisation}\label{sec:MInotation}

The master integrals obtained through IBP reduction are associated with different phase spaces,
depending on the final-state multiplicity of the corresponding antenna. These phase spaces can be
parametrised in terms of the appropriate independent Mandelstam invariants. To simplify the notation
for the corresponding integration measures, we introduce the following shorthand definitions:
\begin{equation}
    \begin{split}
        &s_{12}=s_{12},\quad d\mathcal S_2 = ds_{12}\\
        &s_{123} = s_{12}+s_{13}+s_{23},\quad d\mathcal S_3 = ds_{12}ds_{13}ds_{23}\\
        &s_{1234} = s_{12}+s_{13}+s_{14}+s_{23}+s_{24}+s_{34},\quad d\mathcal S_4 = ds_{12}ds_{13}ds_{14}ds_{23}ds_{24}ds_{34}.
    \end{split}
\end{equation}
Using this notation, the $d$-dimensional differential phase-space measures can be written as
\cite{GehrmannDeRidder0311276,Anastasiou:2002yz},
\begin{equation}\label{eq:phaseSpaces}
    \begin{split}
        &d\Phi_2 = (2\pi)^{2-d}s_{12}^{\frac{d-4}{2}}2^{1-d}d\Omega_{d-1}d\mathcal S_2\delta\!\left(q^2-s_{12}\right)\\
        &d\Phi_3 = (2\pi)^{3-2d}2^{-1-d}\!\left(q^2\right)^{\frac{2-d}{2}}(s_{12}s_{13}s_{23})^{\frac{d-4}{2}}d\Omega_{d-1}d\Omega_{d-2}d\mathcal S_3\delta\!\left(q^2-s_{123}\right)\\
        &\begin{split}
            d\Phi_4 = (2\pi)^{4-3d}\!\left(q^2\right)^{\frac{2-d}{2}}2^{1-2d}&(-\Delta_4)^{\frac{d-5}{2}}\Theta(-\Delta_4)\\ &d\Omega_{d-1}d\Omega_{d-2}d\Omega_{d-3}d\mathcal S_4\delta\!\left(q^2-s_{1234}\right).
         \end{split}
    \end{split}
\end{equation}
For the four-particle phase space, $\Delta_4$ denotes the Gram determinant, which can be expressed in terms of the Källén function as
\cite{GehrmannDeRidder0311276,Anastasiou:2002yz},
\begin{equation}
    \Delta_4=\lambda(s_{12}s_{34}, s_{13}s_{24}, s_{14}s_{23}),\quad \lambda(x, y, z) = x^2+y^2+z^2-2(xy+xz+yz).
\end{equation}
The exponent of $q^2$ in $d\Phi_4$ is fixed by dimensional analysis: a four-particle phase space has mass
dimension $3d-8$, so, with $\Delta_4$ as defined above, the prefactor must be $\left(q^2\right)^{\frac{2-d}{2}}$,
as for $d\Phi_3$. Equation~(A.5) of Ref.~\cite{GehrmannDeRidder0311276} prints the exponent $3-\frac d2$,
which is not consistent with the dimension of $\Delta_4$ as defined here.
The factor $\Theta(-\Delta_4)$, where $\Theta$ is the Heaviside step function, restricts the integration
to the physical region of the four-particle phase space. The angular factors in
Eq.~\eqref{eq:phaseSpaces} are written in terms of $d\Omega_n$, the solid-angle element on the unit
$(n-1)$-sphere, with
\begin{equation}
    \Omega_n=\int d\Omega_n=\frac{2\pi^{\frac n2}}{\Gamma(n/2)}.
\end{equation}

The master integrals are expressed using normalisation factors that collect common
dimensional-regularisation and phase-space factors. We define
\cite{Gehrmann-DeRidder:2004ttg,GehrmannDeRidder0311276},
\begin{equation}
    S_\Gamma = \left(\frac{(4\pi)^\epsilon}{16\pi^2\Gamma(1-\epsilon)}\right)^2.\label{eq:SGamma}
\end{equation}
For phase-space master integrals, it is also convenient to include the integrated two-particle
phase-space volume $P_2$ in the normalisation, defining
\begin{equation}
    S_{\Gamma,2} = P_2 S_\Gamma,\quad P_2 = \int d\Phi_2 = \frac{2^{-3+2\epsilon}\pi^{-1+\epsilon}\Gamma(1-\epsilon)}{\Gamma(2-2\epsilon)}\!\left(q^2\right)^{-\epsilon}.
\end{equation}

\section{$A$-Type Integrals of the $A_2^2$ Set}\label{sec:MIA22}

\subsection{The $A_{22,\text{LO}}$ Integral}

The first master integral required for the $A_2^2$ antenna set is $A_{22,\mathrm{LO}}$, which corresponds
to the product of two massless one-loop bubble integrals. In contrast to the phase-space master
integrals encountered in the integration of real-radiation antennae, $A_{22,\mathrm{LO}}$ is a purely
virtual two-loop integral involving two independent loop momenta. Defining
\begin{equation*}
    q=k_1+k_2,\qquad q^2=s_{12},
\end{equation*}
where $k_1$ and $k_2$ are the momenta of the final-state quark and antiquark, respectively, the
integral is given by \cite{Gehrmann-DeRidder:2005btv},
\begin{equation}
    \begin{split}
        A_{22,\text{LO}} &= \int\frac{d^dp}{(2\pi)^d}\int\frac{d^d\ell}{(2\pi)^d}\frac1{p^2(p-k_1-k_2)^2\ell^2(\ell-k_1-k_2)^2}\\
                         &= -\frac{\mu^{-4\epsilon}}{\left(16\pi^2\right)^2}\Big[B_0\!\left(q^2;0,0\right)\Big]^2\\
                         &= S_\Gamma\!\left(-q^2\right)^{-2\epsilon}\frac{\Gamma^2(1+\epsilon)\Gamma^6(1-\epsilon)}{\Gamma^2(2-2\epsilon)}\!\left(-\frac1{\epsilon^2}\right)\\
                         &= S_\Gamma\!\left(-q^2\right)^{-2\epsilon}\!\left[-\frac1{\epsilon^2}-\frac4\epsilon-12+\mathcal O(\epsilon)\right].
    \end{split}
\end{equation}

Since the propagators depending on $p$ and $\ell$ are independent, the integral factorises into the
product of two identical massless one-loop bubble integrals. It can therefore be written in terms of
the scalar two-point function $B_0(q^2;0,0)$.

\subsection{The $A_3$ Integral}

The $A_3$ master integral is a massless two-loop integral containing three propagators. It is defined
as \cite{Gehrmann-DeRidder:2005btv},
\begin{equation}
    A_3 = \int \frac{d^dp}{(2\pi)^d}\int\frac{d^d\ell}{(2\pi)^d}\frac1{p^2\ell^2(p-\ell-k_1-k_2)^2}.
\end{equation}
Unlike $A_{22,\mathrm{LO}}$, the two loop integrations do not factorise, since the third propagator
depends on both loop momenta. Performing the $\ell$-integration first gives a massless one-loop
bubble, so that,
\begin{equation}
    A_3 = \int \frac{d^dp}{(2\pi)^d}\frac1{p^2}B_0\!\left((p-k_1-k_2)^2;0,0\right).
\end{equation}
Substituting the analytic expression for the massless $B_0$ function and carrying out the remaining
$p$-integration gives
\begin{equation}
    A_3 = S_\Gamma\!\left(-q^2\right)^{1-2\epsilon}\frac{\Gamma(1+2\epsilon)\Gamma^5(1-\epsilon)}{\Gamma(3-3\epsilon)}\!\left(-\frac{1}{2(1-2\epsilon)\epsilon}\right),
\end{equation}
which can be expanded around $\epsilon=0$ as
\begin{equation}
    A_3 = S_\Gamma\!\left(-q^2\right)^{1-2\epsilon}\!\left[-\frac1{4\epsilon}-\frac{13}8+\mathcal O(\epsilon)\right].
\end{equation}

\subsection{The $A_4$ Integral}\label{subsubsec:A4}

The $A_4$ master integral is a massless two-loop integral containing four propagators,
\begin{equation}
    A_4 = \int \frac{d^dp}{(2\pi)^d}\int\frac{d^d\ell}{(2\pi)^d}\frac1{p^2\ell^2(p-k_1-k_2)^2(p-\ell-k_1)^2}.
\end{equation}
As in the $A_3$ integral, the two loop momenta are coupled and the integral does not factorise into a
product of independent one-loop integrals. The $\ell$-integration can, however, be performed first.
The two propagators depending on $\ell$,
\begin{equation}
    \frac1{\ell^2(p-\ell-k_1)^2},
\end{equation}
form a massless one-loop bubble with external momentum $p-k_1$. Hence,
\begin{equation}
    A_4 = \int\frac{d^dp}{(2\pi)^d}\frac1{p^2(p-k_1-k_2)^2}B_0\!\left((p-k_1)^2;0,0\right).
\end{equation}
Substituting the analytic expression for the massless $B_0$ function and evaluating the remaining
$p$-integration gives the following result \cite{Gehrmann-DeRidder:2005btv}, expressed as a Laurent
series in $\epsilon$:
\begin{equation}
    \begin{split}
        A_4 &= S_\Gamma\!\left(-q^2\right)^{-2\epsilon}\frac{\Gamma(1-2\epsilon)\Gamma(1+\epsilon)\Gamma^4(1-\epsilon)\Gamma(1+2\epsilon)}{\Gamma(2-3\epsilon)}\!\left(-\frac1{2(1-2\epsilon)\epsilon^2}\right)\\
            &= S_\Gamma\!\left(-q^2\right)^{-2\epsilon}\!\left[-\frac1{2\epsilon^2}-\frac5{2\epsilon}-\left(\frac{19}2+\frac{\pi^2}6\right)+\mathcal O(\epsilon)\right].
    \end{split}
\end{equation}

\subsection{The $A_6$ Integral}\label{subsubsec:A6}

The final master integral required for the $A_2^2$ antenna set is $A_6$, a massless two-loop integral
containing six propagators. It is defined as \cite{Gehrmann-DeRidder:2005btv},
\begin{equation}
    A_6 = \int\frac{d^dp}{(2\pi)^d}\int\frac{d^d\ell}{(2\pi)^d}\frac1{p^2\ell^2(p-k_1-k_2)^2(p-\ell)^2(p-\ell-k_2)^2(\ell-k_1)^2}.
\end{equation}
In contrast to the preceding master integrals, neither loop integration reduces directly to a scalar
one-loop function of the $B_0$, $C_0$, or $D_0$ type. The integral is the two-loop crossed triangle of
the massless quark form factor. It is evaluated by expressing it as a Mellin--Barnes contour integral,
evaluating the residues, and expanding as a Laurent series around $\epsilon=0$, which gives,
\begin{equation}
    \begin{split}
        A_6 &= \frac{S_\Gamma\left(-q^2\right)^{-2-2\epsilon}}{(2\pi i)^2}\int_{-i\infty}^{i\infty}\int_{-i\infty}^{i\infty}dz_1dz_2\,\frac{\Gamma(-z_1)\Gamma(-z_2)\Gamma(2+2\epsilon+z_1+z_2)\Gamma^2(-\epsilon-z_1)\Gamma^2(-\epsilon-z_2)}{\Gamma(-2\epsilon-z_1-z_2)\Gamma(2+\epsilon+z_1+z_2)}\\
            &= S_\Gamma\!\left(-q^2\right)^{-2-2\epsilon}\!\left[-\frac1{\epsilon^4}+\frac{5\pi^2}{6\epsilon^2}+\frac{27}{\epsilon}\zeta_3+\frac{23\pi^4}{36}+\mathcal O(\epsilon)\right].
    \end{split}
\end{equation}
This expansion agrees with that of the exact expression in $\epsilon$ given in
Ref.~\cite{GehrmannDeRidder0507061}, which involves generalised hypergeometric functions of unit
argument,
\begin{equation}\label{eq:A6exact}
    \begin{split}
        A_6 = S_\Gamma\!\left(-q^2\right)^{-2-2\epsilon}\Bigg[
        &-\frac{\Gamma^3(1-\epsilon)\Gamma(1+\epsilon)\Gamma^4(1-2\epsilon)\Gamma^3(1+2\epsilon)}{\epsilon^4\,\Gamma^2(1-4\epsilon)\Gamma(1+4\epsilon)}\\
        &+\frac{\Gamma^4(1-\epsilon)\Gamma(1+\epsilon)\Gamma(1-2\epsilon)\Gamma(1+2\epsilon)}{2\epsilon^4\,\Gamma(1-3\epsilon)}\,{}_3F_2\!\left(1,-4\epsilon,-2\epsilon;1-3\epsilon,1-2\epsilon;1\right)\\
        &-\frac{4\,\Gamma^4(1-\epsilon)\Gamma(1-2\epsilon)\Gamma(1+2\epsilon)}{\epsilon^2(1+\epsilon)(1+2\epsilon)\,\Gamma(1-4\epsilon)}\,{}_3F_2\!\left(1,1,1+2\epsilon;2+\epsilon,2+2\epsilon;1\right)\\
        &-\frac{\Gamma^5(1-\epsilon)\Gamma(1+2\epsilon)}{2\epsilon^4\,\Gamma(1-3\epsilon)}\,{}_4F_3\!\left(1,1-\epsilon,-4\epsilon,-2\epsilon;1-3\epsilon,1-2\epsilon,1-2\epsilon;1\right)\Bigg].
    \end{split}
\end{equation}

\section{$V$-Type Integrals of the $A_3^1$ Set}\label{sec:MIA31}

The integration of the $A_3^1$ antenna set involves both the one-loop integration and
the integration over the $1\to3$ antenna phase space. After applying the IBP reduction
described in Chapter~\ref{ch:packageframework}, the leading- and subleading-colour
contributions are expressed in terms of the master integrals $V_{5,a}$, $V_{5,b}$, and
$V_8$. These integrals therefore combine loop propagators with the three-particle
phase-space integration and constitute the master-integral basis required for the
real--virtual antenna contributions considered in Subsections~\ref{subsec:A31lead}
and~\ref{subsec:A31sub}. Only the real part of each of them is needed: the one-loop
antennae are built from the real part of the tree--one-loop interference,
$2\operatorname{Re}\!\left(\mathcal M^{0\,\dagger}\mathcal M^1\right)$
(Chapter~\ref{ch:physicsbackground}), so the imaginary parts that the loop integrals
develop in the time-like region $q^2>0$ do not contribute.

\subsection{The $V_{5,a}$ Integral}

The master integral $V_{5,a}$ consists of a massless one-loop bubble integrated over the
three-particle phase space. Denoting the momenta of the final-state quark, antiquark, and
gluon by $k_1$, $k_2$, and $k_3$, respectively, their total momentum is
\begin{equation*}
    q=k_1+k_2+k_3,\qquad q^2=s_{12}+s_{13}+s_{23}.
\end{equation*}
The three final-state momenta are massless, $k_i^2=0$, while $q^2>0$ is the invariant mass
squared of the antenna system. The master integral is defined as
\begin{equation}
    \begin{split}
        V_{5,a} &= \operatorname{Re}\!\left[-i\int d\Phi_3\int\frac{d^dp}{(2\pi)^d}\frac1{p^2(p-k_1-k_2-k_3)^2}\right]\\
                &= \operatorname{Re}\!\left[\int d\Phi_3\,B_0\!\left(q^2;0,0\right)\right].
    \end{split}
\end{equation}
Since the scalar function $B_0\!\left(q^2;0,0\right)$, discussed in
Subsubsection~\ref{subsubsec:B0}, depends only on the fixed total invariant $q^2$, and not
on the individual phase-space invariants, it can be taken outside the three-particle
phase-space integration, which leaves the phase-space volume $R_3=\int d\Phi_3$ of
Eq.~\eqref{eq:R3}. Hence this integral yields \cite{Gehrmann-DeRidder:2005btv},
\begin{equation}
    \begin{split}
        V_{5,a} &= S_{\Gamma,2}\!\left(q^2\right)^{1-2\epsilon}\cos(\pi\epsilon)\!\left(-\frac1\epsilon\right)\frac{\Gamma^6(1-\epsilon)\Gamma(1+\epsilon)}{\Gamma(2-2\epsilon)\Gamma(3-3\epsilon)}\\
                &= S_{\Gamma,2}\!\left(q^2\right)^{1-2\epsilon}\!\left[-\frac1{2\epsilon}-\frac{13}4+\mathcal O(\epsilon)\right].
    \end{split}
\end{equation}

\subsection{The $V_{5,b}$ Integral}

The master integral $V_{5,b}$ has a structure similar to $V_{5,a}$, consisting of a massless
one-loop bubble integrated over the three-particle phase space. The essential difference is
that the external momentum of the bubble is now $k_1+k_3$, rather than the total antenna
momentum $q$. It is therefore defined by,
\begin{equation}
    \begin{split}
        V_{5,b} &= \operatorname{Re}\!\left[-i\int d\Phi_3\int\frac{d^dp}{(2\pi)^d}\frac1{p^2(p-k_1-k_3)^2}\right]\\
                &= \operatorname{Re}\!\left[\int d\Phi_3\,B_0\!\left(s_{13};0,0\right)\right].
    \end{split}
\end{equation}
Unlike $q^2$, which is fixed during the phase-space integration, $s_{13}$ varies over the
three-particle phase space. Consequently, the scalar function $B_0\!\left(s_{13};0,0\right)$
cannot be taken outside the phase-space integral, and $V_{5,b}$ does not factorise into the
product of a scalar one-loop function and the phase-space volume $R_3$.

Using the analytic expression for the massless scalar function $B_0\!\left(s_{13};0,0\right)$
and performing the remaining three-particle phase-space integration gives
\cite{Gehrmann-DeRidder:2005btv},
\begin{equation}
    \begin{split}
        V_{5,b} &= \operatorname{Re}\!\left[\int d\Phi_3\,B_0\!\left(s_{13};0,0\right)\right]\\
                &= S_{\Gamma,2}\!\left(q^2\right)^{1-2\epsilon}\cos(\pi\epsilon)\!\left(-\frac1\epsilon\right)\frac{\Gamma^5(1-\epsilon)\Gamma(1-2\epsilon)\Gamma(1+\epsilon)}{\Gamma(2-2\epsilon)\Gamma(3-4\epsilon)}\\
                &= S_{\Gamma,2}\!\left(q^2\right)^{1-2\epsilon}\!\left[-\frac1{2\epsilon}-4+\mathcal O(\epsilon)\right].
    \end{split}
\end{equation}

\subsection{The $V_8$ Integral}\label{subsubsec:V8}

The third master integral required for the $A_3^1$ antenna set is $V_8$. In contrast to
$V_{5,a}$ and $V_{5,b}$, whose loop integrations correspond to scalar two-point functions,
$V_8$ contains a massless one-loop scalar four-point integral integrated over the
three-particle phase space. It is defined as \cite{Gehrmann-DeRidder:2005btv},
\begin{equation}
    \begin{split}
        V_8 &= \operatorname{Re}\!\left[-i\int d\Phi_3\int\frac{d^dp}{(2\pi)^d}\frac1{p^2(p-k_1)^2(p-k_1-k_3)^2(p-k_1-k_2-k_3)^2}\right]\\
            &= \operatorname{Re}\!\left[-i\int d\Phi_3\,\frac1{s_{12}}\operatorname{Box}(s_{13}, s_{23}, s_{12})\right].
    \end{split}
\end{equation}
The loop integral in this definition is a massless scalar box integral. Using the scalar
one-loop conventions introduced in Subsubsection~\ref{subsubsec:D0}, it may be written in
terms of the $D_0$ function as
\begin{equation}
    V_8 = \left(\frac{(4\pi)^\epsilon\mu^{-2\epsilon}}{16\pi^2}\right)\frac{\Gamma^2(1-\epsilon)\Gamma(1+\epsilon)}{\Gamma(1-2\epsilon)}\int d\Phi_3\,\frac1{s_{12}}\operatorname{Re}\Big[D_0\!\left(0,0,0,q^2;s_{13},s_{23};0,0,0,0\right)\Big].
\end{equation}
Unlike $V_{5,a}$, the scalar box depends explicitly on the phase-space invariants $s_{12}$,
$s_{13}$, and $s_{23}$. The loop integral therefore cannot be factorised from the
three-particle phase-space integration. Performing the remaining phase-space integration and
applying the normalisation introduced above gives \cite{Gehrmann-DeRidder:2005btv},
\begin{equation}
    \begin{split}
        V_8 &= S_{\Gamma,2}\!\left(q^2\right)^{-2-2\epsilon}\Big[-\frac5{2\epsilon^4}+\frac{9\pi^2}{2\epsilon^2}+\frac{89\zeta_3}{\epsilon}+\frac{13\pi^4}{180}+\mathcal O(\epsilon)\Big].
    \end{split}
\end{equation}

\section{$R$-Type Integrals of the Four-Parton Antennae}\label{sec:MIfourParton}

The integration of the tree-level four-parton antennae considered in
Subsections~\ref{subsec:A40}--\ref{subsec:C40} requires integration over the massless $1\to4$
antenna phase space. After mapping the antenna functions onto the corresponding integral family
and applying the IBP reduction described in Chapter~\ref{ch:packageframework}, the integrated
antennae can be expressed in terms of a small set of master integrals, $R_4$, $R_6$,
$R_{8,a}$, and $R_{8,b}$. Unlike the $V$-type master integrals encountered for the
real--virtual $A_3^1$ antennae, these integrals contain no loop integration. They are integrals
over the four-particle phase space, with additional propagator-like factors for $R_6$,
$R_{8,a}$, and $R_{8,b}$. The simplest member, $R_4$, is the four-particle phase-space volume
itself. The following subsections present their definitions and analytic evaluation.

\subsection{The $R_4$ Integral}

The simplest master integral in the four-parton basis is $R_4$, corresponding to the total
massless four-particle phase-space volume,
\begin{equation*}
    R_4=\int d\Phi_4.
\end{equation*}
In $d=4-2\epsilon$ dimensions, the phase-space measure may be expressed in terms of the six
two-particle invariants $s_{ij}=2k_i\cdot k_j$, subject to momentum conservation and the
constraint imposed by the four-particle Gram determinant $\Delta_4$. This gives
\begin{equation}
    R_4 = (2\pi)^{4-3d}\!\left(q^2\right)^{\frac{2-d}{2}}2^{1-2d}\int d\mathcal S_4d\Omega_{d-1}d\Omega_{d-2}d\Omega_{d-3}\,(-\Delta_4)^{\frac{d-5}{2}}\Theta(-\Delta_4)\delta\!\left(q^2-\mathcal S_4\right).
\end{equation}
To evaluate $R_4$, the four-particle phase space is factorised into a two-particle phase space
and a tripole phase-space factor,
\begin{equation*}
    d\Phi_4=d\Phi_2\,d\Phi_T,
\end{equation*}
where $d\Phi_2$ describes the reduced two-particle configuration obtained by clustering the
momenta $p_1,p_3,p_4$ into the mapped radiator momentum $\tilde p_{134}$, with $\tilde p_2$
denoting the corresponding mapped spectator momentum. The factor $d\Phi_T$ is the associated
tripole phase-space measure and contains the remaining degrees of freedom required to resolve
the original four-particle configuration. Following Ref.~\cite{GehrmannDeRidder0311276}, these
degrees of freedom are parametrised by five dimensionless variables with support on the unit
hypercube. We start by defining $y_{ij} = s_{ij}/q^2$ and $y_{ijk} = s_{ijk}/q^2$. Our five
dimensionless variables become \cite{GehrmannDeRidder0311276}:
\begin{itemize}
    \item The rescaled invariant mass of the three-parton cluster, $y_{134} = \frac{s_{134}}{q^2}$;
    \item The longitudinal momentum fraction $z_1 = \frac{p_3\cdot\tilde p_2}{\tilde p_{134}\cdot\tilde p_2}$, which parametrises the fraction associated with parton $3$, measured with respect to the mapped momenta $\tilde p_{134}$ and $\tilde p_2$;
    \item The variable $t$, introduced through the second longitudinal momentum fraction $z_2 = (1-z_1)t = \frac{p_4\cdot\tilde p_2}{\tilde p_{134}\cdot\tilde p_2}$, which parametrises the momentum fraction associated with parton $4$ after the fraction $z_1$ carried by parton $3$ has been specified;
    \item The variable $v$, introduced through the rescaled two-particle invariant $y_{14}=\frac{s_{14}}{q^2} = y_{134}(1-z_1)v$;
    \item The variable $\chi$, which parametrises the allowed range of $y_{13}$, where $y_{13,a}$ and $y_{13,b}$ are the lower and upper boundaries determined by the Gram-determinant constraint as $y_{13} = (y_{13,b}-y_{13,a})\chi+y_{13,a}= \Delta y_{13}\chi+y_{13,a}$.
\end{itemize}
Performing the change of variables from the phase-space invariants to the five dimensionless
variables $y_{134},z_1,t,v,\chi$, including the corresponding Jacobian, maps the physical
integration region onto the unit hypercube $[0,1]^5$ and the master integral becomes
\begin{equation}
    \begin{split}
        R_4 = P_2\frac{16^{-2}\pi^{-5+2\epsilon}}{\Gamma(1-2\epsilon)}\!\left(q^2\right)^{2-2\epsilon}\int_{[0,1]^5}dy_{134}dz_1dtdv&d\chi\,\Big[y_{134}(1-y_{134})(1-z_1)\Big]^{1-2\epsilon}\\
                                                                                                                                                  &\Big[z_1t(1-t)v(1-v)\Big]^{-\epsilon}\Big[\chi(1-\chi)\Big]^{-\frac12-\epsilon}.
    \end{split}
\end{equation}
A useful feature of this parametrisation is that the dependence on the five integration
variables factorises into powers of $x$ and $1-x$. Each integration can therefore be evaluated
using the Euler Beta function,
\begin{equation*}
    \int_0^1 dx\,x^{a-1}(1-x)^{b-1}=B(a,b)=\frac{\Gamma(a)\Gamma(b)}{\Gamma(a+b)}.
\end{equation*}
Performing these integrations and simplifying the resulting Gamma functions gives
\cite{Gehrmann-DeRidder:2005btv,GehrmannDeRidder0311276},
\begin{equation}
    \begin{split}
        R_4 &= P_2\!\left(q^2\right)^{2-2\epsilon}\frac{2^{-11+8\epsilon}\pi^{-\frac72+2\epsilon}\Gamma^3(1-\epsilon)}{\Gamma(3-3\epsilon)\Gamma\!\left(\frac52-2\epsilon\right)}\\
            &= S_{\Gamma,2}\!\left(q^2\right)^{2-2\epsilon}\frac{\Gamma^5(1-\epsilon)\Gamma(2-2\epsilon)}{\Gamma(3-3\epsilon)\Gamma(4-4\epsilon)}\\
            &= S_{\Gamma,2}\!\left(q^2\right)^{2-2\epsilon}\Bigg[\frac1{12}+\frac{59}{72}\epsilon+\mathcal O\!\left(\epsilon^2\right)\Bigg].
    \end{split}
\end{equation}

\subsection{The $R_6$ Integral}

The second master integral required for the four-parton antennae is $R_6$, defined by
\cite{GehrmannDeRidder0311276},
\begin{equation}
    \begin{split}
        R_6 &= \int d\Phi_4\,\frac1{s_{134}s_{234}}\\
            &= (2\pi)^{4-3d}\!\left(q^2\right)^{\frac{2-d}{2}}2^{1-2d}\int d\mathcal S_4d\Omega_{d-1}d\Omega_{d-2}d\Omega_{d-3}\,\frac{(-\Delta_4)^{\frac{d-5}{2}}}{s_{134}s_{234}}\Theta(-\Delta_4)\delta\!\left(q^2-\mathcal S_4\right).
    \end{split}
\end{equation}
Unlike $R_4$, however, $R_6$ does not reduce to a product of elementary Beta-function integrals
under the tripole change of variables introduced above. In particular, $s_{234}$ has a
non-linear dependence on the tripole variables, preventing the factorisation of the integrand
that made the direct evaluation of $R_4$ straightforward. Following
Ref.~\cite{GehrmannDeRidder0311276}, $R_6$ is therefore evaluated indirectly using the optical
theorem.

The optical theorem relates the imaginary part of a forward-scattering integral to the sum over
its possible unitarity cuts. In the present case, the relevant three-loop vacuum-polarisation
integral $I_6$ admits cuts that can be expressed in terms of the four-particle phase-space
integral $R_6$ and the lower-multiplicity contribution involving the two-loop integral $A_4$.
With the conventions of Ref.~\cite{GehrmannDeRidder0311276}, this relation reads
\begin{equation}
    2\operatorname{Im}I_6 = -2P_2\operatorname{Re}A_4 + 2R_6.
\end{equation}
The required expressions for $I_6$ and $A_4$ are known analytically
\cite{GehrmannDeRidder0311276}. In the conventions adopted here, they are
\begin{equation}
    \begin{split}
        I_6 = \frac{(4\pi)^{3\epsilon}}{\left(16\pi^2\right)^3}&\frac{\Gamma^6(1-\epsilon)\Gamma^3(1+\epsilon)}{3\epsilon^3\Gamma^3(2-2\epsilon)}\!\left(-q^2\right)^{-3\epsilon}\\
                                                                      &\Bigg[1+\epsilon+\epsilon^2+\epsilon^3(14\zeta_3-7)+\epsilon^4\left(-67+14\zeta_3+\frac{21\pi^4}{90}\right)+\mathcal O\!\left(\epsilon^5\right)\Bigg]
    \end{split}
\end{equation}
\begin{equation}
    \begin{split}
        A_4 = \left(\frac{(4\pi)^\epsilon}{16\pi^2}\frac{\Gamma(1+\epsilon)\Gamma^2(1-\epsilon)}{\Gamma(1-2\epsilon)}\right)^2&\!\left(-q^2\right)^{-2\epsilon}\Bigg[-\frac1{2\epsilon^2}-\frac5{2\epsilon}-\left(\frac{19}2+\frac{\pi^2}6\right)\\
                                                                                                                              &-\epsilon\left(\frac{65}2+\frac{5\pi^2}6-2\zeta_3\right)-\epsilon^2\left(\frac{211}2+\frac{19\pi^2}6-10\zeta_3\right)+\mathcal O\!\left(\epsilon^3\right)\Bigg].
    \end{split}
\end{equation}
The two-loop master integral $A_4$ was introduced previously in Section~\ref{sec:MIA22}
in the calculation of the two-parton NNLO antennae. Its appearance here illustrates that the
same virtual integral enters the lower-multiplicity cut of the three-loop forward-scattering
integral.

Since the forward-scattering integral is evaluated in the time-like region $q^2>0$, the factors
involving $(-q^2)$ must be analytically continued according to the Feynman $+i0$ prescription.
In the convention of Ref.~\cite{GehrmannDeRidder0311276}, this gives
\begin{equation*}
    \left(-q^2-i0\right)^{-\epsilon}=\left(q^2\right)^{-\epsilon}e^{i\pi\epsilon}.
\end{equation*}
The real and imaginary parts required in the relation above can then be extracted from these
phases. Substituting the real part of $A_4$ and the imaginary part of $I_6$, and denoting by
$\mathcal A_4$ and $\mathcal I_6$ the $\epsilon$-series in square brackets in the expressions
for $A_4$ and $I_6$ above, we solve for $R_6$ and obtain
\begin{equation}
    \begin{split}
        R_6 &= S_{\Gamma,2}\!\left(q^2\right)^{-2\epsilon}\frac{\Gamma^6(1-\epsilon)\Gamma^2(1+\epsilon)}{6}\left(\frac{6\cos(2\pi\epsilon)}{\Gamma^2(1-2\epsilon)}\mathcal A_4+\frac{\Gamma(1-\epsilon)\Gamma(1+\epsilon)\sin(3\pi\epsilon)}{\epsilon^3\pi\Gamma^2(2-2\epsilon)}\mathcal I_6\right)\\
            &= S_{\Gamma,2}\!\left(q^2\right)^{-2\epsilon}\Bigg[-1+\frac{\pi^2}6+\epsilon\!\left(-12+\frac{5\pi^2}6+9\zeta_3\right)+\mathcal O\!\left(\epsilon^2\right)\Bigg].
    \end{split}
\end{equation}

\subsection{The $R_{8,a}$ Integral}\label{subsec:R8aMI}

The master integral $R_{8,a}$ is the first of the two four-particle phase-space master integrals
containing four invariant denominators. It is defined as \cite{GehrmannDeRidder0311276},
\begin{equation}
    \begin{split}
        R_{8,a} &= \int d\Phi_4\,\frac1{s_{13}s_{23}s_{14}s_{24}}\\
                &= (2\pi)^{4-3d}\!\left(q^2\right)^{\frac{2-d}{2}}2^{1-2d}\int d\mathcal S_4d\Omega_{d-1}d\Omega_{d-2}d\Omega_{d-3}\frac{(-\Delta_4)^{\frac{d-5}{2}}}{s_{13}s_{23}s_{14}s_{24}}\Theta(-\Delta_4)\,\delta\!\left(q^2-\mathcal S_4\right)\\
                &   \begin{split}
                        = P_2\!\left(q^2\right)^{-2-2\epsilon}&\frac{16^{-2+\epsilon}\pi^{-5+2\epsilon}}{\Gamma(1-2\epsilon)}\int_{[0,1]^5} dy_{134}dz_1dtdvd\chi\\
                                                              &\Big[(1-y_{134})(1-z_1)\Big]^{-1-2\epsilon}\Big[tvz_1^{-\epsilon}(\Delta y_{13}\chi+y_{13,a})\Big]^{-1}\Big[\chi(1-\chi)\Big]^{-\frac12-\epsilon}.
                    \end{split}
    \end{split}
\end{equation}

In contrast to $R_4$, the resulting integrand does not factorise into independent Beta-function
integrals. Its analytic evaluation instead proceeds through a sequence of integrations involving
hypergeometric functions. Following Ref.~\cite{GehrmannDeRidder0311276}, we perform the
integrations in the order,
\begin{equation*}
    \chi\to v\to y_{134}\to t\to z_1,
\end{equation*}
with an auxiliary integration variable $u\in[0,1]$ introduced at a later stage through an Euler
integral representation of the hypergeometric function.

We begin with the $\chi$ integration. The dependence on $\chi$ occurs through the
Gram-determinant factor and the invariant $y_{13}=y_{13,a}+\Delta y_{13}\chi$, leading to,
\begin{equation}
    I_\chi=\int_0^1 d\chi\,\Big[\chi(1-\chi)\Big]^{-\frac12-\epsilon}(\Delta y_{13}\chi+y_{13,a})^{-1}=\frac{\Delta y_{13}^{-2\epsilon}}{y_{13,a}}\frac{\Gamma^2\!\left(\frac12-\epsilon\right)}{\Gamma(1-2\epsilon)}{}_2F_1\!\left(1, \frac12-\epsilon; 1-2\epsilon; -\frac{\Delta y_{13}}{y_{13,a}}\right)
\end{equation}
where ${}_2F_1$ is the Gauss hypergeometric function. To simplify the argument of the
hypergeometric function, we use the phase-space boundaries,
$y_{13,(a,b)}=y_{134}(A\mp B)$, with
\begin{equation*}
    A=\sqrt{(1-t)(1-v)},\qquad B=\sqrt{z_1tv},\qquad Z=-\frac BA.
\end{equation*}
These definitions imply,
\begin{equation}
    -\frac{\Delta y_{13}}{y_{13,a}} = \frac{4Z}{(1+Z)^2},\quad (1+Z)^2 = \frac{y_{13,a}}{y_{134}A^2}.
\end{equation}
The quadratic transformation of the Gauss hypergeometric function \cite{Erdelyi:1953HTF},
\begin{equation}
    {}_2F_1\!\left(1,b;2b;\frac{4Z}{(1+Z)^2}\right) = (1+Z)^2{}_2F_1\!\left(1,\frac32-b;b+\frac12;Z^2\right)
\end{equation}
together with
\begin{equation}
    16^\epsilon\Delta y_{13}^{-2\epsilon}=\Big[t(1-t)v(1-v)y_{134}^2z_1\Big]^{-\epsilon},
\end{equation}
allows the $\chi$-integrated expression to be written as,
\begin{equation}
    I_\chi = \frac{\pi\Gamma(1-2\epsilon)}{y_{134}\Gamma^2(1-\epsilon)}\!\left(tvy_{134}^2z_1\right)^{-\epsilon}\Big[(1-t)(1-v)\Big]^{-1-\epsilon}{}_2F_1\!\left(1,1+\epsilon;1-\epsilon;\frac{z_1tv}{(1-t)(1-v)}\right).
\end{equation}

Substituting this expression into the phase-space integral eliminates $\chi$ and reduces
$R_{8,a}$ to a four-dimensional integral,
\begin{equation}
    \begin{split}
        R_{8,a} = P_2\!\left(q^2\right)^{-2-2\epsilon}&\frac{16^{-2+\epsilon}\pi^{-4+2\epsilon}}{\Gamma^2(1-\epsilon)}\int_{[0,1]^4}dy_{134}dz_1dtdv\,\Big[y_{134}(1-y_{134})z_1(1-z_1)\Big]^{-1-2\epsilon}\\
                                                      &\Big[t(1-t)v(1-v)\Big]^{-1-\epsilon}{}_2F_1\!\left(1,1+\epsilon;1-\epsilon;\frac{z_1tv}{(1-t)(1-v)}\right).
    \end{split}
\end{equation}

The $v$ integration is performed by expanding the Gauss hypergeometric function in its defining
series and integrating term by term. This yields,
\begin{equation}
    \begin{split}
        I_v &= \int_0^1 dv \Big[v(1-v)\Big]^{-1-\epsilon}{}_2F_1\!\left(1,1+\epsilon;1-\epsilon;\frac{z_1tv}{(1-t)(1-v)}\right)\\
            &= \sum_{n\geq0}\frac{(1+\epsilon)_n}{(1-\epsilon)_n}\!\left(\frac{z_1t}{1-t}\right)^n\frac{\Gamma(-\epsilon-n)\Gamma(-\epsilon+n)}{\Gamma(-2\epsilon)}\\
            &= -\frac{2\Gamma^2(1-\epsilon)}{\epsilon\Gamma(1-2\epsilon)}{}_2F_1\!\left(1,-\epsilon;1-\epsilon;-\frac{z_1t}{1-t}\right).
    \end{split}
\end{equation}

After substituting $I_v$, the $y_{134}$ dependence factorises and can be integrated directly,
leaving a two-dimensional integral over $z_1$ and $t$,
\begin{equation}
    \begin{split}
        R_{8,a} = -P_2\!\left(q^2\right)^{-2-2\epsilon}&\frac{2^{-7+8\epsilon}\pi^{-\frac72+2\epsilon}}{\epsilon^2\Gamma\!\left(\frac12-2\epsilon\right)}\int_{[0,1]^2}dz_1dt\,(1-z_1)^{-1-2\epsilon}\\
                                                       &\Big[z_1t(1-t)\Big]^{-1-\epsilon}{}_2F_1\!\left(1,-\epsilon;1-\epsilon;-\frac{z_1t}{1-t}\right).
    \end{split}
\end{equation}

To perform the remaining integrations, the hypergeometric function is represented through its
Euler integral,
\begin{equation}
    {}_2F_1\!\left(1,-\epsilon;1-\epsilon;-\frac{z_1t}{1-t}\right) = -\epsilon\int_0^1 du\,u^{-1-\epsilon}\!\left[1+\frac{uz_1t}{1-t}\right]^{-1}
\end{equation}
The auxiliary variable $u$ therefore converts the remaining hypergeometric dependence into an
additional ordinary integral over the unit interval. The remaining $t$, $u$, and $z_1$
integrations can then be performed sequentially. After simplifying the resulting Gamma and
generalised hypergeometric functions, one obtains,
\begin{equation}
    \begin{split}
        R_{8,a} = P_2\!\left(q^2\right)^{-2-2\epsilon}&2^{-7+4\epsilon}\pi^{-3+2\epsilon}\csc(\pi\epsilon)\Gamma(-2\epsilon)\Bigg[\frac{2^{8\epsilon}\pi}{\epsilon^2\Gamma^2\!\left(\frac12-2\epsilon\right)}{}_3F_2\!\left(\begin{array}{c} -2\epsilon, -2\epsilon, -2\epsilon \\ 1-2\epsilon, -4\epsilon \end{array};1\right)-\\
                                                      &-\frac{2\Gamma^3(-\epsilon)}{\Gamma(1-4\epsilon)\Gamma(1-\epsilon)\Gamma(1+\epsilon)\Gamma(-3\epsilon)}{}_4F_3\!\left(\begin{array}{c} 1, -\epsilon, -\epsilon, -\epsilon \\ 1-\epsilon, 1+\epsilon, -3\epsilon \end{array};1\right)\Bigg]
    \end{split}
\end{equation}
Rewriting the result in the normalisation used throughout this chapter and expanding about
$\epsilon=0$ gives \cite{GehrmannDeRidder0311276,TM_JoanaReis}, 
\begin{equation}
    \begin{split}
        R_{8,a} = &S_{\Gamma,2}\!\left(q^2\right)^{-2-2\epsilon}\Bigg[\frac6{\epsilon^4}\frac{\Gamma^5(1-\epsilon)\Gamma(1-2\epsilon)}{\Gamma(1-3\epsilon)\Gamma(1-4\epsilon)}{}_4F_3\!\left(\begin{array}{c} 1, -\epsilon, -\epsilon, -\epsilon \\ 1-\epsilon, 1+\epsilon, -3\epsilon \end{array};1\right)-\\
                  &-\frac1{\epsilon^4}\frac{\Gamma^3(1-\epsilon)\Gamma(1+\epsilon)\Gamma^3(1-2\epsilon)}{\Gamma^2(1-4\epsilon)}{}_3F_2\!\left(\begin{array}{c} -2\epsilon, -2\epsilon, -2\epsilon \\ -4\epsilon, 1-2\epsilon \end{array};1\right)\Bigg]\\
                  &=S_{\Gamma,2}\!\left(q^2\right)^{-2-2\epsilon}\Bigg[\frac5{\epsilon^4}-\frac{20\pi^2}{3\epsilon^2}-\frac{126\zeta_3}{\epsilon}+\frac{7\pi^4}{18}+\mathcal O(\epsilon)\Bigg].
    \end{split}
\end{equation}

\subsection{The $R_{8,b}$ Integral}

The second four-denominator master integral is $R_{8,b}$, defined as
\cite{GehrmannDeRidder0311276},
\begin{equation}
    R_{8,b} = \int d\Phi_4\,\frac1{s_{13}s_{134}s_{23}s_{234}}.
\end{equation}
As in the case of $R_6$, the presence of $s_{234}$ prevents the tripole change of variables from
reducing the integral to a convenient factorised form. Following
Ref.~\cite{GehrmannDeRidder0311276}, $R_{8,b}$ is therefore determined indirectly by applying
the optical theorem to the corresponding three-loop vacuum-polarisation integral $I_8$. The
different unitarity cuts of $I_8$ generate contributions involving the two-loop virtual master
integral $A_6$, the real--virtual master integral $V_8$, and the two four-particle phase-space
master integrals $R_{8,a}$ and $R_{8,b}$. With the conventions of
Ref.~\cite{GehrmannDeRidder0311276}, their relation is
\begin{equation}
    2\operatorname{Im}I_8 = -2P_2\operatorname{Re}A_6+4\operatorname{Re}V_8+R_{8,a}+4R_{8,b}.
\end{equation}
The required three-loop integral $I_8$ is known analytically
\cite{GehrmannDeRidder0311276}, while $A_6$ and $V_8$ were introduced previously in
Sections~\ref{sec:MIA22} and~\ref{sec:MIA31}, respectively. Their relevant expressions
are
\begin{equation}
    I_8 = S + \mathcal O(\epsilon^0),\quad \operatorname{Im}S = -\frac{13\pi^4}{90};
\end{equation}
\begin{equation}
    V_8 = -i\int d\Phi_3\,\frac1{s_{12}}\operatorname{Box}(s_{13}, s_{23}, s_{12});
\end{equation}
\begin{equation}
    A_6 = \left(\frac{(4\pi)^\epsilon}{16\pi^2}\frac{\Gamma(1+\epsilon)\Gamma^2(1-\epsilon)}{\Gamma(1-2\epsilon)}\right)^2\!\left(-q^2\right)^{-2-2\epsilon}\Bigg[-\frac1{\epsilon^4}+\frac{5\pi^2}{6\epsilon^2}+\frac{23\zeta_3}\epsilon+\frac{103\pi^4}{180}+\mathcal O(\epsilon)\Bigg].
\end{equation}
Solving the optical-theorem relation for $R_{8,b}$ gives,
\begin{equation}
    R_{8,b}=\frac14\left[2\operatorname{Im}I_8+2P_2\operatorname{Re}A_6-4\operatorname{Re}V_8-R_{8,a}\right].
\end{equation}
Substituting the corresponding $\epsilon$-expansions and rewriting the result in the
$S_{\Gamma,2}$ normalisation used throughout this section yields \cite{GehrmannDeRidder0311276},
\begin{equation}
    R_{8,b} = S_{\Gamma,2}\!\left(q^2\right)^{-2-2\epsilon}\Bigg[\frac3{4\epsilon^4}-\frac{17\pi^2}{12\epsilon^2}-\frac{44\zeta_3}{\epsilon}-\frac{61\pi^4}{60}+\mathcal O(\epsilon)\Bigg].
\end{equation}
